\subsection{The complete level-three spaces in further odd weights}
\label{sec:gl2-general-model}

\begin{proposition}[K17-GL2-GENERAL]\label{prop:gl2-general}
Put $\chi=\chi_{-3}$, and retain the normalized modular forms
\[
 a=\sum_{\alpha\in\mathbf Z[\zeta_3]}q^{N(\alpha)},\qquad
 b=\eta(3\tau)^9/\eta(\tau)^3,\qquad
 D=\eta(\tau)^6\eta(3\tau)^6.
\]
For every integer $e\ge3$, set $d=2e-5$ and $r=\lfloor d/3\rfloor+1$.
Then
\[
 f_j=Da^{d-3j}b^j,\qquad 0\le j<r,
 \tag{G1}
\]
is a basis of $S_{2e+1}(\Gamma_0(3),\chi)$. It is also a basis over
$\mathbf Z$ of the forms in that space whose Fourier coefficients are
integers. If $Q_{e,p}(X)$ is the characteristic polynomial of the
classical $T_p$ on this entire space, define
\begin{align*}
R_{8,7}(X)&=X^2+1570772X-39368833865900,\\
R_{8,13}(X)&=X^2+790365212X+53944448738897860,\\
R_{10,7}(X)&=X^3-279803958X^2-127584561562048500X\\
&\quad+30590253082420358136355000,\\
R_{10,13}(X)&=X^3-8923939638X^2-6797158417228643929140X\\
&\quad+194437243363096745461703318361400,\\
R_{12,7}(X)&=X^3-994032438X^2-313864742351252973300X\\
&\quad-1965630682954936067477202845000,\\
R_{12,13}(X)&=X^3+36771219031962X^2\\
&\quad-143657636498003570201050740X\\
&\quad-7357825592607362422723731960417071505800.
\end{align*}
More precisely,
\[
\begin{split}
 Q_{8,p}&=R_{8,p}^2,\qquad Q_{10,p}=R_{10,p}^2,\\
 Q_{12,7}&=(X+4119710398)R_{12,7}^2,\\
 Q_{12,13}&=(X-41718377201762)R_{12,13}^2.
\end{split}
\tag{G2}
\]
The exact integer matrices, coefficient rows, and evaluations
$Q_{e,p}(1+p^{2e})$ are recorded in
\texttt{hecke\_general/gl2\_general\_certificate.json} and reconstructed by
the accompanying standard-library certificate.
\end{proposition}

\begin{proof}
The primary ring theorem
\cite[paragraph on $\Gamma_1(3)$, p.~12]{K17CleryVanDerGeer2022}
states $M(\Gamma_1(3))=\mathbf C[a,b]$, with generator weights one and
three and the normalizations displayed above. Its odd-weight part is
$M_{\mathrm{odd}}(\Gamma_0(3),\chi)$.
The same source, Remark~15.2, p.~33, identifies $D$ as a weight-six cusp
form with trivial character. To check division by $D$, the eta product
has no interior zeros, begins with $q$, and satisfies
\[
 D\!\left(-\frac1{3\tau}\right)=-27\tau^6D(\tau),
 \qquad D|_6W_3=-D.
\]
It therefore has order exactly one at each of the two cusps.
Both cusp stabilizer translations have lower-right entry one, so
forms of character $\chi$ have integral local orders. Consequently
division by $D$ identifies $S_{2e+1}(\Gamma_0(3),\chi)$ with
$M_{2e-5}(\Gamma_0(3),\chi)$, whose weighted-degree-$d$ monomials are
exactly those in (G1).

The integer expansions used, with their exact normalization, are
\[
\begin{aligned}
 a&=1+6\sum_{n\ge1}\left(\sum_{t\mid n}\chi(t)\right)q^n,\\
 b&=q\prod_{n\ge1}(1-q^{3n})^9(1-q^n)^{-3},\\
 D&=q\prod_{n\ge1}(1-q^n)^6(1-q^{3n})^6.
\end{aligned}
 \tag{G3}
\]
For the first identity, unique factorization in the Eisenstein integers
counts six generators per nonzero ideal; the split, inert, and ramified
prime factors give the divisor sum. The other identities are the
defining eta products. Each $f_j$ has leading term $q^{j+1}$.
Thus the matrix
$L=([q^i]f_j)_{1\le i\le r,\,0\le j<r}$ is integral lower triangular
with diagonal one. Its inverse is integral, and the first $r$
coefficients determine the coordinates of any form in the space.
This proves both completeness and the integral-lattice assertion.

For $p\nmid3$, the classical normalization is
\[
 [q^n](T_pf)=[q^{pn}]f+
 \begin{cases}
 \chi(p)p^{2e}[q^{n/p}]f,&p\mid n,\\
 0,&p\nmid n.
 \end{cases}
 \tag{G4}
\]
See \cite[\S8.2.2, p.~66]{K17RibetStein2011}.
For each column $j$, calculate the first $r$ coefficients of (G4)
from (G3), and solve $Lc_j=v_j$ by the exact recursion
\[
 c_{j,i}=v_{j,i}-\sum_{h<i}L_{ih}c_{j,h}.
 \tag{G5}
\]
Hecke stability of the proved space and invertibility of $L$ make the
result the full operator matrix, not a restriction to selected forms.
Its entries are integers. The finite products truncated after
$q^{pr}$ suffice for this identification. In particular, when $e=12$
and $p=7$, the correction term in (G4) at $n=7$ is included.

For reproducibility, the certificate uses precision $182$, exceeding
all the required $pr$ for $e=6,8,10,12$ and $p=7,13$.
It uses the finite binomial formula for positive product exponents and
$(1-X)^{-s}=\sum_{j\ge0}\binom{s+j-1}{j}X^j$ for negative ones.
It reconstructs every series and matrix, checks (G4) through
$n=\lfloor182/p\rfloor$, and computes the characteristic polynomial by
the exact Faddeev--LeVerrier recursion. Every division in that recursion
is checked, as is the final Cayley--Hamilton zero matrix.
Independent fraction-free determinant calculations check the polynomial
at $X=0,1,1+p^{2e}$. Multiplication of the factors in (G2) gives exactly
these characteristic polynomials. The previously checked weight-thirteen
calculation is reproduced as an additional control. No floating-point,
symbolic-algebra package, or stored Hecke matrix is used by the certificate.
\end{proof}

\paragraph{The exact combined comparison integers.}
Let $P_{e,p}(X)$ denote the characteristic polynomial of the scaled
compact $U(3)$ Hecke operator $p^{e-1}T_p$ on the entire coefficient
space of highest weight $(e-1,0,1-e)$, computed in the independent
unitary certificate. Define
\[
 E_{e,p}=1+p^e+p^{2e},\qquad
 N_{e,p}=P_{e,p}(E_{e,p})Q_{e,p}(1+p^{2e}).
 \tag{G6}
\]
Combining the two exact model certificates gives
\[
\begin{array}{c|l}
e&\gcd(N_{e,7},N_{e,13})\\ \hline
6&2^{15}3^{19}\,7\,13^5\,37\,47^3\\
8&2^{29}3^{31}\,23^3\,401^3\,809\,13687^3\\
10&2^{54}3^{54}\,43\,241^3\,1847\,9496386083^3\\
12&2^{58}3^{71}\,5^2\,7\,13^8\,37\,47\,28447^3\,300455201^3.
\end{array}
\tag{G7}
\]
The standard-library file
\texttt{hecke\_general/hecke\_combined\_certificate.py} reconstructs every
evaluation in (G6) from the full stored characteristic polynomials,
checks agreement with both independently reconstructed model reports,
computes the integer gcds, and verifies each factorization in (G7).
Trial division proves primality of every displayed factor. All factors
for $e=6,8,12$ are less than $2^{31}$. The factor $9496386083$ for
$e=10$ exceeds that range.

\paragraph{Scope and an additional check.}
These are complete finite Hecke models and exact arithmetic identities.
Their application to a Selmer group requires the separate source
admissibility, conductor, transfer, and trace-normalization arguments.
It is enough for that application to cover every admitted representation
pointwise; no multiplicity comparison with the source's larger
\texttt{TrData} is asserted.

At $e=10$, the additional exact classical operator $T_{19}$ is
reconstructed in
\texttt{hecke\_general/gl2\_e10\_p19\_certificate.py}, with coefficient
checks through $q^{20}$ and precision $380$.
Its value $Q_{10,19}(1+19^{20})$ is also divisible by $9496386083$.
Thus this third-prime comparison does not remove the displayed residual
prime. This is an obstruction to this particular numerical unit
criterion, not evidence of a nonzero algebraic $K$-group.
